\documentclass[xcolor=dvipsnames,xcolor=table, 14p]{beamer}

% language support and encoding of input and output
\usepackage[english]{babel}
\usepackage[utf8]{inputenc}
% better font support
\usepackage[T1]{fontenc}
%\usepackage{lmodern}
\usepackage{avant} % URW gothic fonts
% graphical support
\usepackage[table]{xcolor}
\usepackage{graphicx}
\usepackage[table]{xcolor} % beamer already loads this
% typographical improvement
\usepackage{microtype}
\usepackage{tikz}
\usetikzlibrary{shapes, positioning, shapes.arrows}
\usepackage{pdfrender}
\usepackage{anyfontsize}

\definecolor{codegray}{gray}{0.7}
\newcommand{\code}[1]{\colorbox{codegray}{\texttt{#1}}}
\newcommand{\bs}{\textbackslash}

% define colors:
\definecolor{lightblue}{HTML}{7AB8FF}
\definecolor{blue}{HTML}{3F92F2}
\definecolor{green}{HTML}{17BF64}
\definecolor{gray}{HTML}{414444}


% setting up Beamer theme and look:
\usetheme{Madrid}
\setbeamertemplate{items}[circle]
\setbeamertemplate{blocks}[rounded][shadow=true]
\usecolortheme[named=gray]{structure}
\setbeamertemplate{navigation symbols}{}

\setbeamercolor{frametitle}{fg=green}
\setbeamercolor{title}{fg=green}
\setbeamercolor{subtitle}{fg=blue}
\setbeamercolor{author}{fg=blue}

\setbeamerfont{title}{series=\bfseries}
\setbeamerfont{frametitle}{series=\bfseries}
\setbeamerfont{author}{series=\bfseries}

% author and presentation information:
\author[J. M.]{Jiří Moravec}
\title{Escape Excel}
\subtitle{A quick introduction to programming, R, git and reproducibility}
\date{}
\titlegraphic{\includegraphics[width=2cm]{../assets/running_colored.pdf}}

\setbeamertemplate{title page}{%
\begin{tikzpicture}[remember picture,overlay]
% colored rectangle on bottom half of page
\fill[gray]([yshift=15pt]current page.west) rectangle (current page.south east);

% authors name
\node[anchor=east] at ([yshift=25pt]current page.east) (author)
  {\parbox[t]{.8\paperwidth}{\raggedleft%
    \usebeamerfont{author}\usebeamercolor[fg]{author}{\large\insertauthor}}};

% titlegraphics
\node[anchor=south west] 
  at ([yshift=10pt]current page.west) (logo)
  {\parbox[t]{\paperwidth}{\centering%
    \usebeamercolor[fg]{titlegraphic}\inserttitlegraphic}};

% title
\node[anchor=east]
  at ([yshift=-10pt,xshift=-20pt]current page.east) (title)
  {\parbox[t]{0.9\textwidth}{\centering%
    \usebeamerfont{title}\usebeamercolor[fg]{title}\Huge\inserttitle
    }};
\node[anchor=east]
  at ([yshift=-50pt,xshift=-20pt]current page.east) (subtitle)
  {\parbox[t]{.9\paperwidth}{\centering
    \usebeamerfont{subtitle}\usebeamercolor[fg]{subtitle}{\Large \insertsubtitle}}};
\end{tikzpicture}
}



\begin{document}
\section{Introduction}
\begin{frame}[plain]
    \maketitle
\end{frame}


\begin{frame}\frametitle{What is this about?}
    \begin{itemize}
        \setlength\itemsep{1em}
        \item Project organization -- files, folders and data
        \item Introduction to script-based analyses
        \item Introduction to pipelines
        \item Some R and git
    \end{itemize}
\end{frame}


\begin{frame}\frametitle{What is this NOT about?}
    \begin{itemize}
        \setlength\itemsep{1em}
        \item R workshop
        \item Robust tutorial to programming
        \item Extensive tutorial to git\\(most of the time, you are using only 3 commands anyway)
    \end{itemize}
\end{frame}


\begin{frame}{}
    \begin{tikzpicture}[remember picture,overlay]
        \node[anchor=south] at ([yshift=-53pt]current page.center) (running) {
            \includegraphics[height=3cm]{../assets/running.pdf}
            };
        \fill[gray]([yshift=-50pt]current page.west) rectangle (current page.south east);
        \visible<2->{
            \node[anchor=center]
                at ([yshift=50pt]current page.south)
                {\textcolor{blue}{\bf \Large Showing you a direction!}};
            }
        \visible<3->{
            \node[single arrow, draw=gray, anchor=west, fill=gray, shape border rotate=180]
                at (current page.west) (thisway)
                {\textcolor{blue}{\bf \Large This way!}};
            \node[right = 2cm of thisway] (start) {};
            }

    \end{tikzpicture}
\end{frame}


\begin{frame}<1-2>[label=content]\frametitle{Program:}
    \begin{enumerate}
        \setlength\itemsep{1em}
        \item \alert<2>{Motivation}
        \item \alert<3>{Project organization}
        \item \alert<4>{Pipelines and script-based analyses}
        \item \alert<5>{Introduction to programming}
        \item \alert<6>{git}
    \end{enumerate}
\end{frame}
\section{Motivation}

\begin{frame}\frametitle{Motivation:}
    \begin{tikzpicture}[align=center, node distance = 2cm, remember picture,overlay]
        \node[anchor=north] at (current page.center) (excel) {\includegraphics[height=1.5cm]{../assets/Excel_logo.pdf}};
        \node[above left = of excel] (data) {\Large data};
        \node[above = of excel] (code) {\Large code};
        \node[above right = of excel] (results) {\Large results};
        \draw[->, line width=3pt, gray]
            (data) edge (excel)
            (code) edge (excel)
            (results) edge (excel);
        \visible<2->{
            \node at (current page.center) (bad) {\textcolor{red}{\bf \fontsize{50}{60}\selectfont BAD}};
            }
    \end{tikzpicture}
\end{frame}

\begin{frame}\frametitle{Motivation:}
    \begin{center}
    \Large \only<3->{Collaborator\_}analysis\only<2->{\_final}\only<4->{Final}.xls
    \end{center}
\end{frame}

\begin{frame}<1-2>[label=motivation_folders]\frametitle{
    \only<1-2>{Motivation:}
    \only<3>{Project organization}
    }
    \begin{columns}
        \centering
        \begin{column}{0.33\textwidth}
            \begin{itemize}
                \setbeamertemplate{itemize items}{
                    \includegraphics[height=0.5cm]{../assets/folder.pdf}
                    }
                \item Data
                \item Scripts
                \item Results
                \only<2->{
                \item Intermediaries
                \item Documentation
                }
            \end{itemize}
        \end{column}
        \begin{column}{0.33\textwidth}
            \begin{itemize}
                \setbeamertemplate{itemize items}{
                    \includegraphics[height=0.7cm]{../assets/file2.pdf}
                    }
                \item makefile
                \item run.sh
                \item README.md
            \end{itemize}
        \end{column}
    \end{columns}
\end{frame}


\section{Project organization}

%\begin{frame}\frametitle{Program:}
%    \begin{enumerate}
%        \setlength\itemsep{1em}
%        \item Motivation
%        \item \alert{Project organization}
%        \item Pipelines and script-based analyses
%        \item Introduction to programming
%        \item git
%    \end{enumerate}
%\end{frame}

\againframe<3>{content}

\begin{frame}\frametitle{Project organization}
    Make sure you have:
    \begin{enumerate}
        \item a single folder for entire project
        \item reasonable naming scheme
        \item good folder structure
    \end{enumerate}
\end{frame}

\againframe<3>{motivation_folders}

\begin{frame}\frametitle{Project organization}
    \begin{tabular}{ll}From  & \it A Guide to Reproducible Code\\& \it in Ecology and Evolution\end{tabular}\uncover<2->{\textcolor{red}{\bf Notice the names!}}
    \begin{columns}
        \centering
        \begin{column}{0.20\textwidth}
            \begin{itemize}
                \setbeamertemplate{itemize items}{
                    \includegraphics[height=0.5cm]{../assets/folder.pdf}
                    }
                \item data
                \item doc
                \item figs
                \item output
                \item R
            \end{itemize}
        \end{column}
        \begin{column}{0.5\textwidth}
            \begin{itemize}
                \setbeamertemplate{itemize items}{
                    \includegraphics[height=0.7cm]{../assets/file2.pdf}
                    }
                \item 01\_download-data.R
                \item 02\_clean-data.R
                \item 03\_exploratory-analysis.R
                \item 04\_fit-models.R
                \item 05\_generate-figures.R
                \item README.md
            \end{itemize}
        \end{column}
    \end{columns}
\end{frame}

\againframe<4>{content}

\begin{frame}\frametitle{Pipelines}
    \begin{tikzpicture}[remember picture,overlay]
    \tikzstyle{edge}=[->, line width=3pt, gray]
        \node[anchor=west]
            at ([yshift=-2.5cm, xshift=0.25cm]current page.north west)
            (data) {\Large data};
        \node[right = 2cm of data] (results) {\Large results};
        \draw[edge] (data) edge node[above] {analysis} (results);
        \uncover<2->{
            \node[below = of results] (data2) {\Large data};
            \node[right = 2cm of data2] (results2) {\Large results};
            \draw[edge] (data2) edge node[above] {analysis} (results2);
            \draw[edge] (results) edge (data2);
            }
        \uncover<3->{
        \node[below = of results2] (data3) {\Large data};
        \node[right = 2cm of data3] (results3) {\Large results};
        \draw[edge] (data3) edge node[above] {analysis} (results3);
        \draw[edge] (results2) edge (data3);
        }
        \uncover<4->{
            \node at ([yshift=1cm]current page.south) {\Large Functions, \alert<5->{Scripts}, smaller or bigger analyses}; 
            }
    \end{tikzpicture}
\end{frame}

\begin{frame}\frametitle{Script-based analyses}
\begin{itemize}
    \item repeatability
    \item documented process
    \item (including any data modification and filtering)
    \item smaller effort to modify
    \item can be repeated by just running the run script
\end{itemize}
\end{frame}

\againframe<5>{content}

\begin{frame}\frametitle{Introduction to programming}
Why \textcolor{blue}{\bf R}?
\begin{itemize}
    \item easy to get data into R
    \item closer to statistics than other programming languages
    \item graphical capabilities
    \item huge library of packages for essentially anything
    \item FREE and OPEN SOURCE
\end{itemize}
\end{frame}

\begin{frame}\frametitle{Introduction to programming}
Helpful tips:
\begin{itemize}
    \item Scripts:
    \begin{itemize}
        \item self-contained
        \item doing one thing well rather than many things badly
    \end{itemize}
    \item Functions:
    \begin{itemize}
        \item self-contained
        \item small
        \item use them instead of repeating code
        \item use them for a well-defined tasks and concepts
        \item documentation and comments
    \end{itemize}
    \item Data:
    \begin{itemize}
        \item text human readable files
        \item not binary proprietary formats
        \item modify with scripts
        \item never overwrite
    \end{itemize}
\end{itemize}
\end{frame}

\begin{frame}\frametitle{Introduction to programming}
Helpful tips (cont):
\begin{itemize}
    \item Paths:
    \begin{itemize}
        \item use relative rather than absolute
        \item relative \code{data\bs my\_data.csv}
        \item absolute: \code{c:\bs my\_projects\bs data\bs my\_data.csv}
    \end{itemize}
    \item Figures:
    \begin{itemize}
        \item generate with scripts
        \item easier to modify
        \item easier to change once data or analysis changes
    \end{itemize}
\end{itemize}
\end{frame}



\begin{frame}{}
    \begin{center}
        \textcolor{blue}{\LARGE \bf Live coding excercise!}
    \end{center}
\end{frame}

\againframe<6>{content}

\begin{frame}\frametitle{Git}
\begin{columns}
    \begin{column}{0.45\textwidth}
        Git is:
        \begin{itemize}
            \item version control system
            \item track all changes in text files (data, code)
            \item tracks where the changes were done
            \item FREE and OPEN-SOURCE
        \end{itemize}
    \end{column}
    \begin{column}{0.45\textwidth}
        \centering
        \includegraphics[width=0.9\textwidth]{../assets/git_logo.eps}
    \end{column}
\end{columns}
\end{frame}

\begin{frame}\frametitle{Git}
Git commands:
\begin{itemize}
    \item \code{git init}
    \item \code{git clone <path or url>}
    \item \code{git add file1 file2} or \code{git add -A}
    \item \code{git commit -m "Information about commit"}
    \item \code{git push remote\_repository branch}
    \item \code{git pull remote\_repository branch}
    \item \code{git status}
\end{itemize}
\end{frame}

\begin{frame}\frametitle{Sources and guides}
{\Large \bf R}
\begin{itemize}
    \item Start with:  \url{http://tryr.codeschool.com}
    \item Continue with: \url{http://tryr.codeschool.com}
    \item Advanced R: \url{http://adv-r.had.co.nz/}
\end{itemize}
{\Large \bf git}
\begin{itemize}
    \item Start with: \url{https://guides.github.com/introduction/git-handbook/}
    \item Documentation: \url{https://git-scm.com/docs}
    \item Book: \url{https://git-scm.com/book/en/v2}
\end{itemize}
{\large \bf A Guide to Reproducible Code in Ecology and Evolution}
\begin{itemize}
    \item \url{https://tinyurl.com/reproducibleCode}
    \item \href{https://www.britishecologicalsociety.org/wp-content/uploads/2017/12/guide-to-reproducible-code.pdf?fbclid=IwAR00qqChWaXk1o1_gjOtPkcD_y21n3BtLItdM07B8VyjmdrOLAHC6UEk0aQ}{original link}
\end{itemize}
\end{frame}

\begin{frame}{}
    \begin{center}
        \textcolor{blue}{\LARGE \bf I want your feedback!}
    \end{center}
\end{frame}
\end{document}
